\section{Exact barriers and honest unresolved inputs}
\label{sec:barriers}

A general framework is useful only if it distinguishes formalization problems from
missing mathematics and from impossible interfaces.  This section records the most
important boundaries.  None of them says that the associated theorem is unprovable in
\(\iPRA\); each says that a particular generic route cannot supply the required data.

\subsection{Cardinality does not uniformly produce a cutoff}

Suppose \(P_e(n)\) is a decidable family with at most one solution.  It is tempting to
expect a uniform algorithm that, from \(e\) and the cardinality bound, returns \(B(e)\)
such that every solution is below \(B(e)\).  This is impossible.

\begin{barrier}[No computable cardinality-to-cutoff operator]
\label{barrier:cardinality-cutoff}
There is no total computable operator which, given a code for a decidable predicate
\(P(n)\) together with a valid certificate that \(P\) has at most one solution, always
returns a bound \(B\) such that
\[
  P(n)\quad\Longrightarrow\quad n<B.
\]
In particular, no primitive-recursive operator has this property.
\end{barrier}

\begin{proof}
Let \(M_e\) be the Turing machine with code \(e\).  Define the decidable predicate
\[
  P_e(n)\;:\Longleftrightarrow\;
  M_e\text{ halts for the first time after exactly }n\text{ steps}.
\]
Bounded simulation makes \(P_e\) primitive recursive in \((e,n)\).  It has at most one
solution, and the proof of that fact is uniform: a deterministic machine cannot have
two distinct first halting times.

Apply the hypothetical operator to obtain \(B(e)\).  Simulate \(M_e\) for \(B(e)\)
steps.  If it halts, answer yes.  If it does not halt, it can never halt later, because
a later first halting time \(n\ge B(e)\) would be a solution of \(P_e\) outside the
promised cutoff.  This decides the halting problem, a contradiction.
\end{proof}

Thus a finiteness or cardinality theorem is not enough for complete enumeration.  A
height bound, search cutoff, or complete-list certificate is genuinely additional
data.  This is the precise issue in generic forms of Siegel, Faltings, and
Mordell--Weil enumeration.

\subsection{Exact zero tests from approximate names}

\begin{barrier}[No finite-use exact zero test]
\label{barrier:zero-test}
There is no extensional finite-use procedure that, from an arbitrary regular Cauchy
name of a real number known only to be dyadic rational, always returns the exact bit
\([x=0]\).
\end{barrier}

\begin{proof}
Run the supposed procedure on a regular name for zero.  Because the output is a finite
computation, it reads only rows through some horizon \(H\).  Choose a positive dyadic
\(\delta=2^{-M}\) with \(M\) so large that a regular name for \(\delta\) may output
zero on every row through \(H\) and switch to sufficiently accurate approximations to
\(\delta\) later.  For example, take \(M>H+2\); the discrepancy from zero at the early
rows is below their allowed error.  The two names are indistinguishable to the
procedure, so it returns the same bit, although one value is zero and the other is not.
Contradiction.
\end{proof}

Promises about the abstract value do not repair the information interface.  Saying
``the represented number happens to be algebraic'' or ``dyadic'' is weaker than
providing a finite exact presentation from which equality is decidable.  Exact Puiseux
leading exponents, multiplicities, branch ranks, and similar discrete outputs can
encode such a zero test.  Therefore an unconditional exact-output theorem for arbitrary
Cauchy coefficients is not available.  One must require exact algebraic splitting
data or return a relational approximate object.

\subsection{No universal primitive-recursive evaluator}

The diagonal proof in \cref{barrier:no-universal-pr} blocks another common shortcut.
A code for an arbitrary primitive-recursive function is not automatically executable
by a total primitive-recursive universal semantics.  Static specialization, finite
query tables, or a restricted fuel-bounded interpreter are the admissible alternatives.

\subsection{Metastability does not imply a tail rate}

Finite-window stability is compatible with arbitrarily late future oscillations.  A
uniform conversion from every metastability bound to a Cauchy modulus would therefore
supply information not present in the premise.  In applications, a metastable theorem
is often exactly what a finite consumer needs, and it should be retained as such.  A
limit name requires a separate tail modulus, summable motion estimate, monotonicity
argument, or other global input.

\subsection{Negative nonzeroness does not supply positive separation}

The guarded reciprocal theorem shows what division needs: a rational lower bound
\(\lambda\).  The proposition \(x\ne0\) alone gives no algorithm for finding
\(\lambda\) from an arbitrary Cauchy name.  This is another discontinuity at zero.
Accordingly, every inverse, normalization, pivot, or branch-isolation interface lists
its apartness data explicitly.

\subsection{Qualitative compactness does not supply a net}

A theorem that a space is compact in a strong ambient theory may have no
primitive-recursive finite-net producer with a modulus usable in \(\iPRA\).  The
\(\FRNF\) compactness interface is conditional on such data.  For boxes, simplices,
finite complexes, rational polytopes, and many explicitly presented compact spaces,
the data are easy to construct.  For an arbitrary quotient or function space, they may
be the main mathematical theorem.

\subsection{Deep analytic and arithmetic producers}

Finite contour integration, exact rational linear algebra, and bounded enumeration can
formalize the consequences of a zero-free region, exceptional-zero separation, height
bound, or covering list.  They do not manufacture those inputs.  In particular:

\begin{itemize}
\item stronger prime-number-theorem rates require correspondingly strong zero-free
regions and growth estimates;
\item Siegel--Walfisz and Bombieri--Vinogradov uniformity require substantive control of
exceptional characters/zeros, not merely better bookkeeping;
\item general Chebotarev needs effective field algorithms and analytic estimates with
field dependence tracked;
\item general Siegel and Faltings enumeration needs an effective height producer; and
\item general central-limit and stable-law statements need tail, smoothing,
quantization, or classification inputs beyond finite-law arithmetic.
\end{itemize}

\subsection{Effectivity obstacles are not independence proofs}

Failure of a generic producer does not prove that no \(\iPRA\) proof of a particular
closed theorem exists.  A theorem may have a special argument, a nonuniform constant,
or a certificate discovered by external mathematics and then verified in \(\iPRA\).
The barriers above justify exact claim wording:

\begin{quote}
``The finite verifier and conditional implication are formalized; the required
producer is not supplied by this framework.''
\end{quote}

They do not justify the stronger wording ``the theorem is unprovable in \(\iPRA\)''
without a separate proof-theoretic argument.
